% GENERATED FILE. Edit canonical lessons, then run npm run generate:books.
% canonical-source-hash: fnv1a64:825030a825f34067
% canonical-lessons: HI-C310-bichurna, HI-C310-bhatakna, HI-C310-talas, HI-C310-salamat, HI-C310-dilasa

\chapter{Lost and Found}
\label{ch:hi-lost-and-found}
\hlchaptermodality{310}

\begin{chapteropening}
\textbf{By the end of this chapter:} \emph{I can tell how a family was separated and found again.}

Complete the last lesson of chapter 310: I can tell how a family was separated and found again.
\end{chapteropening}

\section[\texorpdfstring{\hi{बिछुड़ना} (bichuṛnā)}{बिछुड़ना (bichuṛnā)}]{\hi{बिछुड़ना} (bichuṛnā) — to be separated}

\label{lesson:HI-C310-bichurna}



\begin{quote}
Before the new one: say the Hindi for a victim, then the Hindi for refuge.
\end{quote}

\subsection*{You'll want to know: \hi{बिछुड़ना}}
\textbf{\hi{बिछुड़ना}} — \emph{bichuṛnā} — \textquotedblleft{}to be separated\textquotedblright{}.

\begin{quote}
\hi{मेले} \hi{में} \hi{एक} \hi{बच्चा} \hi{अपनी} \hi{माँ} \hi{से} \hi{बिछुड़} \hi{गया।}
\end{quote}

\begin{quote}
\emph{mele meṁ ek baccā apnī māṁ se bichuṛ gayā.}
\end{quote}

\begin{quote}
At the fair a child got separated from his mother.
\end{quote}

Past (he …): \hi{वह} \hi{बच्चा} \hi{बिछुड़} \hi{गया}, \emph{vah baccā bichuṛ gayā}.

\begin{cousinweb}
\textbf{\hi{बिछुड़ना}} (\emph{bichuṛnā}) is inherited from \emph{*vikṣuṭati}, a form CDIAL reconstructs (CDIAL 11651). The one you are parted from takes \textbf{\hi{से}}: \textbf{\hi{माँ} \hi{से} \hi{बिछुड़} \hi{गया}}.
\end{cousinweb}

\begin{grammarlens}[title={What you've built}]
One more word for telling what happened.
\end{grammarlens}

\subsection*{Your turn}
\begin{itemize}
\raggedright
  \item \emph{Say it:} \emph{bichuṛnā}
  \item \emph{Say it:} \emph{bichuṛnā}, once more
  \item \emph{Say it:} \emph{mele meṁ ek baccā apnī māṁ se bichuṛ gayā.} — At the fair a child got separated from his mother.
  \item \emph{Recall:} say the Hindi for a victim, then the Hindi for refuge, then say \emph{bichuṛnā} again
\end{itemize}

\begin{tcolorbox}[breakable,colback=teal!4,colframe=teal!35!black,title={Before you move on}]
What does \hi{बिछुड़ना} mean? (\textquotedblleft{}To be separated\textquotedblright{}.) Say it once more.
\end{tcolorbox}
\section[\texorpdfstring{\hi{भटकना} (bhaṭaknā)}{भटकना (bhaṭaknā)}]{\hi{भटकना} (bhaṭaknā) — to wander; to lose one's way}

\label{lesson:HI-C310-bhatakna}



\begin{quote}
Before the new one: say the Hindi for refuge, then the Hindi for to be separated.
\end{quote}

\subsection*{You'll want to know: \hi{भटकना}}
\textbf{\hi{भटकना}} — \emph{bhaṭaknā} — \textquotedblleft{}to wander; to lose one's way\textquotedblright{}.

\begin{quote}
\hi{बच्चा} \hi{भटककर} \hi{नदी} \hi{के} \hi{किनारे} \hi{पहुँचा।}
\end{quote}

\begin{quote}
\emph{baccā bhaṭakkar nadī ke kināre pahũcā.}
\end{quote}

\begin{quote}
The child wandered off and ended up on the riverbank.
\end{quote}

Past (he …): \hi{वह} \hi{रास्ता} \hi{भटक} \hi{गया}, \emph{vah rāstā bhaṭak gayā}.

\begin{cousinweb}
\textbf{\hi{भटकना}} (\emph{bhaṭaknā}) is set by CDIAL under Sanskrit \emph{bhraṣṭa}, \textquotedblleft{}fallen, ruined\textquotedblright{} (CDIAL 9655). It is wandering without knowing the way: \textbf{\hi{रास्ता} \hi{भटकना}}, to lose one's way.
\end{cousinweb}

\begin{grammarlens}[title={What you've built}]
One more word for telling what happened.
\end{grammarlens}

\subsection*{Your turn}
\begin{itemize}
\raggedright
  \item \emph{Say it:} \emph{bhaṭaknā}
  \item \emph{Say it:} \emph{bhaṭaknā}, once more
  \item \emph{Say it:} \emph{baccā bhaṭakkar nadī ke kināre pahũcā.} — The child wandered off and ended up on the riverbank.
  \item \emph{Recall:} say the Hindi for refuge, then the Hindi for to be separated, then say \emph{bhaṭaknā} again
\end{itemize}

\begin{tcolorbox}[breakable,colback=teal!4,colframe=teal!35!black,title={Before you move on}]
What does \hi{भटकना} mean? (\textquotedblleft{}To wander; to lose one's way\textquotedblright{}.) Say it once more.
\end{tcolorbox}
\section[\texorpdfstring{\hi{तलाश} (talāś)}{तलाश (talāś)}]{\hi{तलाश} (talāś) — a search}

\label{lesson:HI-C310-talas}



\begin{quote}
Before the new one: say the Hindi for to be separated, then the Hindi for to wander.
\end{quote}

\subsection*{You'll want to know: \hi{तलाश}}
\textbf{\hi{तलाश}} — \emph{talāś} — \textquotedblleft{}a search\textquotedblright{}.

\begin{quote}
\hi{दो} \hi{दिन} \hi{तक} \hi{बच्चे} \hi{की} \hi{तलाश} \hi{चली।}
\end{quote}

\begin{quote}
\emph{do din tak bacce kī talāś calī.}
\end{quote}

\begin{quote}
The search for the child went on for two days.
\end{quote}

Feminine.

\begin{cousinweb}
\textbf{\hi{तलाश}} (\emph{talāś}) is usually traced to Persian \emph{talāš}, \textquotedblleft{}a search, an effort\textquotedblright{}, which may itself come from Turkic. \textbf{\hi{की} \hi{तलाश} \hi{में}} is in search of.
\end{cousinweb}

\begin{grammarlens}[title={What you've built}]
One more word for telling what happened.
\end{grammarlens}

\subsection*{Your turn}
\begin{itemize}
\raggedright
  \item \emph{Say it:} \emph{talāś}
  \item \emph{Say it:} \emph{talāś}, once more
  \item \emph{Say it:} \emph{do din tak bacce kī talāś calī.} — The search for the child went on for two days.
  \item \emph{Recall:} say the Hindi for to be separated, then the Hindi for to wander, then say \emph{talāś} again
\end{itemize}

\begin{tcolorbox}[breakable,colback=teal!4,colframe=teal!35!black,title={Before you move on}]
What does \hi{तलाश} mean? (\textquotedblleft{}A search\textquotedblright{}.) Say it once more.
\end{tcolorbox}
\section[\texorpdfstring{\hi{सलामत} (salāmat)}{सलामत (salāmat)}]{\hi{सलामत} (salāmat) — safe and sound}

\label{lesson:HI-C310-salamat}



\begin{quote}
Before the new one: say the Hindi for to wander, then the Hindi for a search.
\end{quote}

\subsection*{You'll want to know: \hi{सलामत}}
\textbf{\hi{सलामत}} — \emph{salāmat} — \textquotedblleft{}safe and sound\textquotedblright{}.

\begin{quote}
\hi{आख़िर} \hi{में} \hi{बच्चा} \hi{सलामत} \hi{मिला।}
\end{quote}

\begin{quote}
\emph{ākhir meṁ baccā salāmat milā.}
\end{quote}

\begin{quote}
In the end the child was found safe and sound.
\end{quote}

\begin{cousinweb}
\textbf{\hi{सलामत}} (\emph{salāmat}) is usually traced, through Persian, to Arabic \emph{salāma}, \textquotedblleft{}safety\textquotedblright{}. It describes a person who has come through unhurt.
\end{cousinweb}

\begin{grammarlens}[title={What you've built}]
One more word for telling what happened.
\end{grammarlens}

\subsection*{Your turn}
\begin{itemize}
\raggedright
  \item \emph{Say it:} \emph{salāmat}
  \item \emph{Say it:} \emph{salāmat}, once more
  \item \emph{Say it:} \emph{ākhir meṁ baccā salāmat milā.} — In the end the child was found safe and sound.
  \item \emph{Recall:} say the Hindi for to wander, then the Hindi for a search, then say \emph{salāmat} again
\end{itemize}

\begin{tcolorbox}[breakable,colback=teal!4,colframe=teal!35!black,title={Before you move on}]
What does \hi{सलामत} mean? (\textquotedblleft{}Safe and sound\textquotedblright{}.) Say it once more.
\end{tcolorbox}
\section[\texorpdfstring{\hi{दिलासा} (dilāsā)}{दिलासा (dilāsā)}]{\hi{दिलासा} (dilāsā) — comfort; consolation}

\label{lesson:HI-C310-dilasa}



\begin{quote}
Before the new one: say the Hindi for a search, then the Hindi for safe and sound.
\end{quote}

\subsection*{You'll want to know: \hi{दिलासा}}
\textbf{\hi{दिलासा}} — \emph{dilāsā} — \textquotedblleft{}comfort; consolation\textquotedblright{}.

\begin{quote}
\hi{पड़ोसियों} \hi{ने} \hi{माँ} \hi{को} \hi{दिलासा} \hi{दिया।}
\end{quote}

\begin{quote}
\emph{paṛosiyoṁ ne māṁ ko dilāsā diyā.}
\end{quote}

\begin{quote}
The neighbours comforted the mother.
\end{quote}

Masculine.

\begin{cousinweb}
\textbf{\hi{दिलासा}} (\emph{dilāsā}) is usually traced to Persian \emph{dilāsā}, \textquotedblleft{}a consoling\textquotedblright{}, whose first part is \emph{dil}, the heart, as in \textbf{\hi{दिल}}. \textbf{\hi{दिलासा} \hi{देना}} is to comfort someone.
\end{cousinweb}

\begin{grammarlens}[title={What you've built}]
One more word for telling what happened.
\end{grammarlens}

\subsection*{Your turn}
\begin{itemize}
\raggedright
  \item \emph{Say it:} \emph{dilāsā}
  \item \emph{Say it:} \emph{dilāsā}, once more
  \item \emph{Say it:} \emph{paṛosiyoṁ ne māṁ ko dilāsā diyā.} — The neighbours comforted the mother.
  \item \emph{Recall:} say the Hindi for a search, then the Hindi for safe and sound, then say \emph{dilāsā} again
\end{itemize}

\begin{tcolorbox}[breakable,colback=teal!4,colframe=teal!35!black,title={Before you move on}]
What does \hi{दिलासा} mean? (\textquotedblleft{}Comfort; consolation\textquotedblright{}.) Say it once more.
\end{tcolorbox}
